% !PRINTED 244
\weberpage{244}
%<<CONT
worin $c_{h,k}$ die zu bestimmenden Koeffizienten sind, die nach 3. der Bedingung
\[
c_{h,k} = c_{k,h}
\]
genügen. Um sie zu bestimmen, setzen wir in (29)
\[
x = u, \quad y = v,
\]
wodurch $F_p$ verschwindet, und erhalten
\weq{30}{(u^p - v)(u - v^p) = \sum^{h,k} c_{h,k} u^h v^k.}
Hieraus folgt zunächst, daß
\[
c_{p,p} = 0
\]
sein muß, da nach (28) bei der Entwickelung (der linken Seite) nach Potenzen von $q$ die Potenz $q^{-2(p^2+p)}$ nicht vorkommen kann, und wir können (29) jetzt in die Form setzen:
\weq{31}{(u^p - v)(u - v^p) = \sum_{0,p}^{h} \sum_{0,h-1}^{k} c_{h,k}(u^h v^k + u^k v^h) + \sum_{0,p-1}^{h} c_{h,h} u^h v^h.}
Hierin sind die aus (25) und (27) sich ergebenden Entwickelungen von
\[
u^h v^k + u^k v^h, \quad u^h v^h
\]
nach Potenzen von $q$ einzusetzen, deren Anfangsglieder
\[
q^{-2(hp+k)}, \quad q^{-2h(p+1)}
\]
die Koeffizienten 1 haben.
Auf der rechten Seite von (31) kommen nicht zwei Glieder vor, deren Entwickelung mit derselben Potenz von $q$ beginnt, denn aus
\[
hp + k = h'p + k'
\]
folgt $k \equiv k'\ (\mathrm{mod}\ p)$ und mithin, da $k$ und $k' < p$ sind, $k = k'$, $h = h'$.
Ordnet man daher die Reihen, welche die beiden Seiten von (31) darstellen, nach aufsteigenden Potenzen von $q$, und setzt dann die Koeffizienten gleich hoher Potenzen einander gleich, so erhält man eine Reihe linearer Gleichungen für die Unbekannten $c_{h,k}$, von denen jede folgende nur \gesperrt{eine} neue Unbekannte enthält, und diese mit dem Koeffizienten 1. Die aus der linken Seite sich ergebenden bekannten Glieder dieser Gleichungen sind nach (28) lauter \gesperrt{durch $p$ teilbare ganze Zahlen}, und es ergeben sich also für die $c_{h,k}$ ebenfalls ganzzahlige, durch $p$ teilbare Zahlwerte. Demnach haben wir
\weq{32}{F_p(x, y) = (x^p - y)(x - y^p) - p \sum_{0,p}^{h,k} a_{h,k} x^h y^k,}
%CONT>>
